(1) Bose 分布为 $f(\varepsilon)=[e^{(\varepsilon-\mu)/(k_BT)}-1]^{-1}$，基态能量取零，因此平衡化学势 $\mu\le0$。三维动量态密度积分给激发态数
\[
N_{\rm ex}(T,z)=\frac{V}{\lambda_T^3}g_{3/2}(z),\quad
g_s(z)=\sum_{\ell=1}^\infty\frac{z^\ell}{\ell^s},\quad z=e^{\mu/(k_BT)}\le1.
\]
当 $\mu\to0^-$，$g_{3/2}(z)\to\zeta(3/2)$，故激发态饱和数
\[
N_{\rm ex}^{\max}=\frac V{\lambda_T^3}\zeta(3/2).
\]
临界点 $N=N_{\rm ex}^{\max}(T_c)$，解得
\[
T_c=\frac{2\pi\hbar^2}{mk_B}
\left(\frac{n}{\zeta(3/2)}\right)^{2/3}.
\]

(2) 在 $T<T_c$ 的热力学极限，化学势趋于 $\mu=0$，超出激发态容量的粒子占据 $\boldsymbol k=0$ 基态，因此
\[
\frac{N_0}{N}=1-\frac{V\zeta(3/2)}{N\lambda_T^3}
=1-\left(\frac{T}{T_c}\right)^{3/2}.
\]
该宏观占据由 Bose 对称化与三维态密度共同产生；本模型中粒子无相互作用，故不能解释成分子间吸引造成的凝聚。

(3) 激发态内能是
\[
U=\frac32\frac{Vk_BT}{\lambda_T^3}\zeta(5/2)
=\frac32Nk_BT\frac{\zeta(5/2)}{\zeta(3/2)}
\left(\frac{T}{T_c}\right)^{3/2}.
\]
基态能量为零，压力由自由气体动量流或热力学关系给出
\[
P=\frac{2U}{3V}=\frac{k_BT}{\lambda_T^3}\zeta(5/2).
\]
固定 $N,V$ 下 $T_c$ 是常数，故
\[
C_V=\left(\frac{\partial U}{\partial T}\right)_{N,V}
=\frac{15}{4}Nk_B\frac{\zeta(5/2)}{\zeta(3/2)}
\left(\frac{T}{T_c}\right)^{3/2}.
\]
不能给凝聚粒子额外附加经典 $3k_BT/2$ 动能，否则将重复计数。

(4) $T\to0^+$ 时 $N_0/N\to1$、$U,P,C_V\to0$；当 $T\uparrow T_c$，$N_0/N\to0$，
\[
C_V(T_c^-)=\frac{15}{4}Nk_B\frac{\zeta(5/2)}{\zeta(3/2)}.
\]
热力学极限中的零化学势与尖锐相变，在有限体积中被平滑；粒子相互作用会改变激发谱、化学势和临界温度，因此这些理想公式只适用于明确给定的无相互作用模型。
